% !TeX root = ../main.tex
\chapter{Causal perturbations and extension of distributions}\label{f:extension}

Send in two weak sources and ask for the part of a detector reading that
vanishes when either source is switched off. For a classical nonlinear
wave this question requires a retarded inverse. For a quantum wave it
matters whether the reading is a change of the mean field or a correlation
of two insertions. The same retarded tree will occur in the change of
mean. A pair of quantum vertices, even when each vertex has already been
defined, will require a different operation.

\section{Differentiating a retarded solution}
A formal series in $\lambda,s,r$ with coefficients in an algebra $A$ is
a list of coefficients $a_{ijk}\in A$ for the monomials $\lambda^i s^j
r^k$, with nonnegative indices. In a product the coefficient at $(i,j,k)$
is the sum over decompositions of that triple between the factors, in the
specified order if $A$ is noncommutative. This is a finite sum even when
the series do not converge anywhere. Differentiation and setting the
parameters to zero extract coefficients, with the usual factorials:
$\partial_s\partial_r|_0$ extracts the $sr$ coefficient, whereas
$\partial_s^2|_0$ gives twice the $s^2$ coefficient (leaving $\lambda$
formal in both cases).

For a differentiable functional $F$ of a smooth field its $k$th variation,
when it exists, is
\[
 F^{(k)}[\phi](h_1,\ldots,h_k)=
 \left.\partial_{s_1}\cdots\partial_{s_k}
 F[\phi+s_1h_1+\cdots+s_kh_k]\right|_{s_1=\cdots=s_k=0} .
\]
The directions $h_i$ are smooth; compact supports or switching functions
keep all displayed integrals finite. For the local polynomial
$F[\phi]=\int\chi\phi^3\,\mathrm{d}V/3!$, the first variation is
$\tfrac12\int\chi\phi^2 h_1\,\mathrm{d}V$ and the second is
$\int\chi\phi h_1 h_2\,\mathrm{d}V$. The latter has kernel
$\chi(x)\phi(x)\delta(x,y)$, where the delta is relative to spacetime
volume. The two variations act at the same point because $F$ is local,
not because the functions $h_1,h_2$ were chosen to be identical.

Take $P=\Box_g-m^2-\xi R$ in a smooth globally hyperbolic spacetime
without boundary.
Use zero incoming data and smooth compactly supported $\chi,f,h$; an
initial slice before the supports suffices. Add
$-\lambda F-\int j\phi\,\mathrm{d}V$ to the free action with bulk
quadratic form $\tfrac12\int\phi P\phi\,\mathrm{d}V$. The equation and its
retarded form are
\begin{equation}
 P\phi=j+\tfrac\lambda2\chi\phi^2,\qquad
 \phi=G_Rj+\tfrac\lambda2 G_R(\chi\phi^2),\qquad j=sf+rh .
 \label{f:ext-retarded}
\end{equation}
Here $PG_R=1$ on tests. The background at $j=0$ is zero. For a nonzero
background solution $\bar\phi$ at source $j_0$, subtract its equation:
the linear operator on $\eta=\phi-\bar\phi$ is $P-\lambda\chi\bar\phi$,
and the new source is $j-j_0$. Its retarded inverse must then be used,
with fixed initial and exterior conditions. We keep the zero background
so that powers of $\lambda$ count only cubic vertices.

Write $u_f=G_R f$, $u_h=G_R h$. Substituting the leading solution on the
right of \eqref{f:ext-retarded} gives
\[
 \phi=su_f+ru_h+\tfrac\lambda2 G_R\!\left(\chi(su_f+ru_h)^2\right)
          +O(\lambda^2).
\]
The $\lambda sr$ coefficient is $G_R(\chi u_f u_h)$, not half this
quantity: there are two assignments of source labels to the square. A
smooth detector $k$ of compact support measures this coefficient as
\begin{equation}
 \int k\,G_R(\chi u_f u_h)\,\mathrm{d}V
       =\int (G_A k)\chi u_f u_h\,\mathrm{d}V .
 \label{f:ext-tree}
\end{equation}
This uses formal self-adjointness and retarded--advanced adjointness on
the same domain. It locates the signal at an overlap of the futures of
both pulses and the past of the detector. Drawing a line for each inverse
gives a rooted tree: two inputs meet, and their product propagates to the
output. Iterating substitutes a new branch at a field factor; it never
pairs two previously existing factors. No state has entered.

\begin{exercise}
In a flat massless region take $\chi=f=h=1$ in the portion of the causal
past of a detector between the initial slice and elapsed time $T$. Compute
the $\lambda sr$ and $\lambda s^2$ field coefficients there. The data
outside that domain may be cut off smoothly.
\end{exercise}
\noindent\emph{Hint.} In the spatially constant calculation $P=-\partial_t^2$
and $u_f=u_h=-t^2/2$. Propagate $t^4/4$ once more with kernel $-(T-t)$;
the mixed answer is $-T^6/120$, the unmixed one half as large.

\section{Ordered insertions: the tree and the two-line loop}
Choose a centered quasifree Hadamard state with two-point function $W$ for
the free field $\Phi$, and set $\hbar=1$. First use Wick ordering in this
reference state. The first subtraction is explicit:
$:\!\Phi(x)\Phi(y)\!:=\Phi(x)\Phi(y)-W(x,y)\mathbf1$.
In matrix elements between vectors generated from the quasifree cyclic
vector by finitely many smeared fields, the ordered pairing rule derived on
p.~\pageref{f:wick-pairings} now pairs each
unsmeared field only with a smeared one. Such a contraction is smooth:
a singular direction cannot have a zero covector in the smeared slot by
\eqref{f:wf-spectrum}. These matrix elements can be restricted to $x=y$;
the result is the Wick square as a distribution in the remaining point.
For higher Wick powers the colons likewise exclude all internal pairings.
Local covariant Wick powers subtract a geometric parametrix rather than
the full $W$; the resulting finite local normalizations must be specified
as well \cite{HWick}.
For a nonderivative potential the Hamiltonian insertion has the opposite
sign to the added Lagrangian term. Thus the finite-time Dyson convention
is $U=\mathcal T\exp[-i\int (j\Phi+\lambda\chi:\!\Phi^3\!:/3!)\,\mathrm{d}V]$,
with later times on the left. Here and below a field at a point is kernel
notation for a spacetime-smeared insertion, not a point operator.

For a change of the mean, the observable evolves as $U^*\Phi(x)U$ up to
the observation time. To first order in an insertion its change is $i$
times the ordered integral of its commutator with $\Phi(x)$. Since
$i[\Phi(y),\Phi(x)]=E(x,y)\mathbf1$, a linear source alone shifts the
field by $u=G_Rj$; further source commutators vanish. Write $U=U_j U_{\rm rel}$
with $U_j$ generated by the source alone. The cubic insertion in the
relative evolution is conjugated by $U_j$ up to its insertion time, so
its field is $\Phi(y)+u(y)$. Its commutator is then
\[
 i[:\!(\Phi(y)+u(y))^3\!:,\Phi(x)]
       =3E(x,y):\!(\Phi(y)+u(y))^2\!: .
\]
Retarded ordering supplies $G_R(x,y)$. Of the terms on the right only
$u(y)^2$ is quadratic in the source. The $3/3!$ from the commutator and
the $2$ from its $sr$ coefficient reproduce \eqref{f:ext-tree}. This
calculation explains how a retarded line arises from quantum evolution;
it is not a prescription to contract two fields with $G_R$.

For a time-ordered correlation the single contraction is instead
\begin{equation}
 G_F=W+iG_A=W^T+iG_R,\qquad P_xG_F(x,y)=i\delta(x,y) .
 \label{f:ext-feynman}
\end{equation}
Indeed $W-W^T=i(G_R-G_A)$, so the first expression equals $W$ to the
future of $y$ and the second equals $W^T$ to its past, off coincidence.
At spacelike separation they agree. This choice defines the free scalar
contraction on the diagonal as well, without multiplying steps at zero;
there is \emph{no} extra $i$ multiplying the time-ordered expectation.
The delta follows by applying $P_x$ to the formula, not to $W$ alone.

For bosonic local insertions the multilinear time-ordered kernels $T_n$
obey \emph{causal factorization}. Split the labels into two nonempty sets
$I,J$. On supports for which no point of group $I$ lies in the causal
past of a point of group $J$,
\begin{equation}
 T_n(F_1,\ldots,F_n)=T_{|I|}(F_I)\,T_{|J|}(F_J) .
 \label{f:ext-factor}
\end{equation}
The right side is the ordinary quantum product, not another ordering
across the groups. To see the reach of this prescription, take a tuple
not all at the same point. The finitely many distinct points have a
maximal element in their causal order. Put \emph{all} labels at this
element into $I$, the others into $J$. Closedness of the causal relation
in a globally hyperbolic spacetime gives open neighborhoods with the
same separation. Such neighborhoods cover the complement of the total
diagonal. Lower products already include their partial coincidences;
causal factorization at lower orders and commutativity at spacelike
separation give agreement on overlaps. This is the off-diagonal step of
the Epstein--Glaser induction \cite{BF,HWT,FR}. A spatial division does
not order the two groups in time.

In the product of two Wick cubes, contract $j$ fields at $x$ with $j$
fields at $y$. The choices of fields and their bijections number
$\binom3j^2j!$, respectively $1,9,18,6$ for $j=0,1,2,3$. Each cross
contraction contributes $W(x,y)$ in the indicated operator order, hence
$G_F(x,y)$ after time ordering away from coincidence. In particular the
two-line term of the quadratic Dyson coefficient is
\begin{equation}
 -\frac{\lambda^2}{4}\int \chi(x)\chi(y)\,G_F(x,y)^2
       :\!\Phi(x)\Phi(y)\!:\,\mathrm{d}V_x\,\mathrm{d}V_y ,
 \qquad x\ne y\ \text{initially} .
 \label{f:ext-fish}
\end{equation}
The factor is $18(-i\lambda)^2/[2!(3!)^2]$; the remaining fields have
not been paired. Two lines between the same two vertices make a loop.
It is not the power of $\lambda$ that identifies the loop: at order
$\lambda^2$ the classical recursion gives the tree
$\tfrac12 G_R(\chi u\,G_R(\chi u^2))$, with $u=G_Rj$.
At a null pair off the diagonal $G_F$ locally selects one of $W,W^T$, so its square
passes the product test of Chapter~\ref{f:microlocal}. At coincidence
\eqref{f:ext-feynman} implies that $\operatorname{WF}(G_F)$ contains all
conormals of the diagonal: differentiation cannot create directions and
$\delta(x,y)$ has them all. The product test therefore fails there. The
object to be constructed is an extension of the \emph{already defined}
separated-vertex distribution, not a point value of $G_F^2$.

\begin{exercise}
Replace $\chi$ in the cubic interaction by $a\chi_1+b\chi_2$. Find the
coefficient of $\lambda^2 ab$ in the two-line term. If an extension of
$G_F^2$ is changed by $c\delta(x,y)$, what change does this coefficient
undergo? Can its nonlocal part persist for disjoint supports?
\end{exercise}
\noindent\emph{Solution.} Exchange $x,y$ in one cross term. The two give
$-\tfrac12$ times the integral in \eqref{f:ext-fish} with switchings
$\chi_1(x)\chi_2(y)$ and no $\lambda^2$. The local change is
$-\tfrac c2\int\chi_1\chi_2:\!\Phi^2\!:\,\mathrm{d}V$. Only this latter
term must vanish for disjoint supports.

\section{Normal scaling and the missing Taylor moments}
At two coincident bulk points in dimension $D$, take coordinates
$(X,z)=(y,x-y)$ in a chart. The diagonal is $z=0$; approaching it means
scaling the $D$ relative coordinates, not scaling $X$. We first fix the
model normal space $\mathbb R^k$. Let $t^\circ$ be a distribution on its
punctured version. For a test $\psi$ supported off zero set
$\psi_a(z)=a^{-k}\psi(z/a)$ and $t^\circ_a(\psi)=t^\circ(\psi_a)$.
Pairing a regular function this way gives the function $t^\circ(az)$;
the factor $a^{-k}$ removes the volume Jacobian.

\begin{definition}
The \emph{scaling degree} at zero is
\[
 \operatorname{sd}_0(t^\circ)=\inf\{d\in\mathbb R:
      a^d t^\circ_a\longrightarrow0\ \text{in distributions as }a\downarrow0\}.
\]
For a distribution defined at zero as well the definition uses all tests,
not just punctured-space tests. A nonzero homogeneous distribution of
degree $-b$ has scaling degree $b$; logarithmic factors in its scaling
law do not increase the infimum.
\end{definition}

\begin{theorem}[Point extension with finite scaling degree, Brunetti--Fredenhagen \cite{BF}]
\label{f:ext-theorem}
A distribution on a punctured neighborhood in $\mathbb R^k$ of finite
scaling degree $\rho$ at zero has an extension of the same scaling degree.
With this degree requirement it is unique for $\rho<k$. For $\rho\ge k$
the freedom is exactly addition of
\begin{equation}
 \sum_{|\alpha|\le\lfloor\rho-k\rfloor}c_\alpha\,\partial^\alpha\delta_0 .
 \label{f:ext-jets}
\end{equation}
No symmetries or field identities have yet been imposed on the constants.
\end{theorem}

Two extensions agree on all tests supported off zero. By the structure
theorem for point-supported distributions \cite{Schwartz,HormI}, their
difference is a finite sum of delta derivatives. On a scaled test a term
$\partial^\alpha\delta_0$ gives
$(-1)^{|\alpha|}a^{-k-|\alpha|}\partial^\alpha\psi(0)$. Prescribing a
Taylor jet of $\psi$ can isolate any highest derivative. This proves the
ambiguity bound and explains why merely requiring \emph{some} extension
would not imply uniqueness even below threshold.

Here is how to use the existence statement. If $m=\lfloor\rho-k\rfloor
\ge0$, take a compactly supported smooth $w$ equal to one near zero and
subtract the jet from the test:
\[
 \psi^\flat(z)=\psi(z)-w(z)\sum_{|\alpha|\le m}
                   \frac{z^\alpha}{\alpha!}\partial^\alpha\psi(0) .
\]
On an annulus of radius $a$ this gains $m+1$ powers of $a$. To see why
that suffices, localize the annulus with a fixed smooth annular cutoff
$\beta(z/a)$. Its pairing with $t^\circ$ is
$a^k t^\circ_a(\beta(\cdot)\psi^\flat(a\,\cdot))$. For any $d>\rho$,
$a^d t^\circ_a$ is bounded on tests of bounded seminorms on that fixed
annulus; Taylor's formula bounds each needed derivative of
$\psi^\flat(a\,\cdot)$ there by $O(a^{m+1})$. The shell is therefore
$O(a^{k+m+1-d})$. Choose $\rho<d<k+m+1$; the bound is summable on a
smooth annular partition with radii $2^{-j}$. If $\rho<k$, no jet has to
be subtracted and the same argument works with exponent $k-d>0$.
Assigning values to the finitely many subtracted Taylor moments completes
the extension; the theorem controls its scaling on all tests. It is the
codimension $k$, not the dimension of the whole configuration space, that
pays for the gain from the volume element.

\section{A marginal subtraction and its scale residue}
In the flat massless four-dimensional theory the time-ordered contraction
is $(4\pi^2)^{-1}(z^2+i0)^{-1}$ with $z^2=|\mathbf z|^2-(z^0)^2$.
Its square on punctured relative space has scaling degree four; it is
still a boundary value \emph{on} the noncoincident light cone. To perform
a subtraction by an ordinary radial integral, first take a different,
Euclidean distribution with the same marginal degree:
\[
 t_E^\circ(z)=\frac{1}{16\pi^4|z|^4},\qquad z\in\mathbb R^4\setminus\{0\},
\]
where $|z|$ is Euclidean length. Its integral against a test constant near
zero would contain $\int_0^1\mathrm{d}r/r$. The theorem calls for one
Taylor moment. For a smooth radial cutoff profile $w$ equal to one near
zero and zero at large radius, and a length $\ell>0$, put
\begin{equation}
 R_\ell(\psi)=\frac1{16\pi^4}\int_{\mathbb R^4}
  \frac{\psi(z)-\psi(0)w(|z|/\ell)}{|z|^4}\,\mathrm{d}^4z .
 \label{f:ext-subtraction}
\end{equation}
The numerator is $O(|z|)$ near zero by the mean value bound. Thus the
radial integral is absolutely convergent, controlled by a first derivative
seminorm of the test there and its supremum away from zero. This defines
a distribution and leaves every punctured-space pairing unchanged. The
subtraction is needed for existence; choosing which constant to assign
to the removed moment is a separate finite choice.

Only the cutoff changes when $\ell$ changes. Since the area of the unit
three-sphere is $2\pi^2$, radial integration gives
\begin{align}
 (R_{\ell'}-R_\ell)(\psi)
 &=-\frac{\psi(0)}{8\pi^2}\int_0^\infty
       [w(r/\ell')-w(r/\ell)]\,\frac{\mathrm{d}r}{r}
   =-\frac{\log(\ell'/\ell)}{8\pi^2}\psi(0) .
 \label{f:ext-residue}
\end{align}
For the last equality differentiate the integral in $\log\ell'$: after
$u=r/\ell'$ the derivative is $-\int_0^\infty w'(u)\,\mathrm{d}u=1$.
The integral is zero at $\ell'=\ell$. Changing variables $z=a v$ in
\eqref{f:ext-subtraction} now yields the full scaling law
\[
 a^4(R_\ell)_a=R_{\ell/a}=R_\ell+\frac{\log a}{8\pi^2}\delta_0 .
\]
Multiplication by any additional positive power of $a$ gives zero in the
limit, and off zero the scaling degree was already four. Thus the
extension preserves scaling degree but not exact homogeneity. Adding a
finite multiple of $\delta_0$ cannot cancel the logarithm because the
delta itself scales homogeneously. The residue just computed uses a
Euclidean radial measure; it is not the complex scale change of the
Lorentzian time-ordered contraction.

\begin{exercise}
Recover the Lorentzian scale change without using a Euclidean sphere on a
light cone. Write $q=z^2$, $q_+=q+i0$, and extend the bracket below
without increasing its scaling degree. Why does the result
\[
 L_\ell=-\frac1{64\pi^4}\Box\!\left[\frac{\log(q_+/\ell^2)}{q_+}\right]
\]
agree off zero with the scalar kernel in \eqref{f:ext-fish}? Determine
$L_{\ell'}-L_\ell$.
\end{exercise}
\noindent\emph{Solution.} Off the cone use
$\Box h(q)=4q h''(q)+8h'(q)$: for $h(q)=\log(q/\ell^2)/q$ the result
is $-4/q^2$. At any nonzero point of the cone $\mathrm{d}q\ne0$, so the
one-variable upper boundary value in $q$ gives the same identity there.
The bracket has scaling degree two (with a logarithm); extend it uniquely
at the tip before differentiating. Two derivatives raise scaling degree
by at most two, and the punctured square already has degree four.
Finally \eqref{f:ext-feynman} gives $\Box(1/q_+)=4\pi^2 i\delta_0$.
The difference is $i\log(\ell'/\ell)\delta_0/(8\pi^2)$. The factor $i$
and its sign depend on the Lorentzian source convention.

\section{Where a contact can live, and which identities it must obey}
For the total diagonal of $n$ bulk vertices in spacetime dimension $D$,
there are $k=(n-1)D$ relative coordinates and $D$ coordinates along the
diagonal. Near a smooth collision stratum use a product chart $(X,z)$ with
$z=0$ on the stratum. Dilate only $z$, using tests
$\psi_a(X,z)=a^{-k}\psi(X,z/a)$. With a finite normal scaling bound
$\rho$ uniform on compact tangential test families, two extensions
obeying that bound can differ locally only by terms of the form
\eqref{f:ext-jets} in $z$, with distributional coefficients in $X$.
This follows from the local structure of distributions supported on a
submanifold and the same scaled-jet test. Smoothness or a prescribed
wavefront set for the coefficients is an additional condition. If they are
singular at the edge of a face, that smaller stratum needs separate
analysis. Power counting classifies possible contacts once the distributions
away from the stratum and their normal estimates are known.

If all $n$ bulk vertices collide \emph{on} a face of codimension $q$,
the common point must move to the face as well: $k=(n-1)D+q$. Vertices
constrained from the start to that face instead have total diagonal of
codimension $(n-1)(D-q)$ in their own configuration space. For example,
draw a plane in the flat four-dimensional fish calculation. The uncut
kernel has no dependence on the normal position $X_\perp$ of the common
point. Under joint scaling of $(z,X_\perp)$ it still has degree four,
whereas $\delta^{(4)}(z)\delta(X_\perp)$ has degree five. With that bound
there is no extra sheet mass term. More generally an artificial division
supplies no new physical coupling: split a switching function smoothly and
sum its terms using the same whole-space extension. A true wall with its
own action, boundary condition or fields supplies new local data. Its
face and corner counterterms must be tested against those data, not
inferred by declaring the wall transparent. A bulk existence theorem does
not solve a theory on an arbitrary boundary.

Symmetries test the \emph{extended} products. At zero interaction the
massless flat scalar has the shift current $j^\mu=-\partial^\mu\Phi$,
with $j^0=\partial_t\Phi$. It is conserved in the bulk field equation;
its time-ordered insertion with a field is not divergence-free at
coincidence. With the derivative normalization of the free contraction,
\begin{equation}
 \partial_\mu^x\langle\mathcal T j^\mu(x)\Phi(y)\rangle
   =-\Box_x G_F(x,y)=-i\delta(x,y) .
 \label{f:ext-ward}
\end{equation}
The first derivative of the ordering has no contact proportional to the
equal-time field--field commutator, which vanishes. The divergence then
detects $[j^0(t,\mathbf x),\Phi(t,\mathbf y)]=-i\delta^{(3)}(\mathbf
x-\mathbf y)$. A test equal to one near $y$ reads the coefficient $-i$.
Thus extending a classically conserved expression by setting its ordered
divergence to zero would give the wrong generator on the inserted field.
The cubic interaction explicitly breaks this shift; such breaking is not
a quantum anomaly. More generally, locality, covariance, power counting,
unitarity and Ward identities restrict the finite terms allowed in an
extension \cite{HWick,HWT,FR}.

In a gauge theory the obstruction has the cochain form of
Chapter~\ref{f:gauge-complex}. Let $s$ be the nilpotent linearized
symmetry differential on admissible local densities, raising ghost number
by one. Suppose all Ward identities through order $n-1$ hold. The first
possible breaking $\mathcal A_n$ is a local contact of ghost number one:
applying the symmetry consistency identity a second time leaves
$s\mathcal A_n=0$ at this order, modulo permitted total derivatives,
since the terms containing lower-order breakings vanish. Adding a finite
local counterterm $B$ of ghost number zero at order $n$ changes the
breaking to $\mathcal A_n+sB$ in this sign convention. Thus the relevant
cohomology is closed local breakings modulo such exact ones, not the
space of arbitrary distributions on the diagonal.

For example take an Abelian potential $A$ and odd ghost $c$ with
$sA_\mu=\partial_\mu c$, $sc=0$. Suppose the dimension and other
normalizations permit the counterterm $B=\tfrac\kappa2\int A_\mu A^\mu\,\mathrm{d}V$.
Then
\[
 sB=\kappa\int A^\mu\partial_\mu c\,\mathrm{d}V
   =-\kappa\int c\,\nabla_\mu A^\mu\,\mathrm{d}V
      +\kappa\int_{\partial M} c A^\mu\,\mathrm{d}\Sigma_\mu ,
\]
with the directed surface element in the last integral. A candidate
breaking $+\kappa\int c\,\nabla_\mu A^\mu\,\mathrm{d}V$ is removable in
the compactly supported bulk problem. On a true wall the flux must vanish
by the physical conditions or be canceled in the combined boundary
identity before the same total derivative can be discarded. The example
identifies a coboundary under stated assumptions; it neither produces an
anomaly nor excludes other cohomology classes.

Finally all these equations are statements in a specified theory and, if
several global sectors exist, in a specified sector. Even if two sectors
have indistinguishable local singularities, no step in the Taylor
subtraction assigns a relative probability or phase to them. To combine
sector answers one needs a common observable and an independently chosen
measure: mixing normalized states needs nonnegative weights summing to
one; adding amplitudes is a different prescription. Conservation and the
extension theorem dictate neither. Equality after a weighted sum would
also be weaker than equality in each sector, since two discrepancies can
cancel.

\begin{exercise}
As a power-counting test consider $|z|^{-6}$ off zero in Euclidean four
space. Which normal contacts can change a degree-preserving extension?
If the result must also be a rotation-invariant scalar, which survive?
Does rotation invariance determine their coefficients?
\end{exercise}
\noindent\emph{Solution.} The excess degree is $6-4=2$. Before symmetry,
delta derivatives through order two are allowed. There is no invariant
vector for the order-one coefficient and an invariant symmetric
rank-two tensor is proportional to the Euclidean metric. Thus
$c_0\delta_0+c_2\Delta\delta_0$ remains. Neither constant is fixed by
this count, which imposed no scaling law beyond the degree bound.
